%! @@TITLE@@ — Unit @@UNITINT@@, Lesson @@LESSONID@@ — Group Activity
% GRADUAL RELEASE: the "you do together" block — 16 minutes in groups of 3-4, followed by an
% 8-minute whole-class debrief that the Debrief box at the end captures.
% ONE activity for the entire class. There are NO tiers, no differentiated versions, and no
% `tierbox` — differentiation happens in how the teacher circulates, not on the page.
% The vocabulary was GIVEN in the guided notes, so it is used freely here: the work is
% APPLYING it to a context students have not seen, and defending every answer from the data.
% Scenario 2 carries the CRUX — the item that surfaces the target misconception.
% All displays are PRE-DRAWN: never ask a student to sketch, draw, or construct a graph.
%
% PAGE LOCKSTEP with activity_key/: \answerspace{H}{} reserves exactly H in the blank and
% prints the red answer in the key, so the two files paginate identically by construction.
% Sizing: 1.0cm ~ 2 handwritten lines, 1.5cm ~ 3, 2.0cm ~ 4. Keep the macro definition, every
% \answerspace height, and every \boxguard byte-identical in both files.
\documentclass[10pt]{article}
\usepackage{@@PREFIX@@-article}
\usepackage{@@PREFIX@@-boxes}
\newcommand{\answerspace}[2]{\par\nopagebreak\noindent\begin{minipage}[t][#1][t]{\linewidth}%
  \color{keyred}\bfseries #2\end{minipage}\par}

\begin{document}

\pageheader{Unit @@UNITINT@@, Lesson @@LESSONID@@}{Group Activity: TODO Activity Title}
% NAMESTRIP: no \namedateperiod here — the name row lives on the cover only.

\vspace{0.10in}

\boxguard[10]
\begin{headlinebox}{sky}
\small
\textbf{TODO Activity Title.} TODO: one motivating data context, 2--3 sentences, ending with
the ground rule --- what groups must be ready to show before they write an answer down.
\end{headlinebox}

\vspace{0.10in}

\boxguard[28]
\begin{scenariobox}[1. TODO Scenario One]{navy}
\small
TODO: the pre-drawn table or display. Scenario 1 is deliberately the clean case, so every
group has something correct on the page before they hit the crux.

\vspace{6pt}
\begin{enumerate}[label=\textbf{\alph*.}, leftmargin=1.8em, itemsep=5pt]
  \item TODO: identification, using the notes' vocabulary. \blank{3.0cm}
  \item TODO: a value to read or compute from the display.
  \item TODO: an answer that must be justified \emph{from the data}.
        \answerspace{1.5cm}{}
\end{enumerate}
\end{scenariobox}

\vspace{0.10in}

\boxguard[26]
\begin{scenariobox}[2. TODO Scenario Two --- the crux]{navy}
\small
TODO: the contrast case, on the same context.

\vspace{6pt}
\begin{enumerate}[label=\textbf{\alph*.}, leftmargin=1.8em, itemsep=5pt]
  \item TODO: the setup that sets the trap.
  \item \textbf{TODO: THE CRUX.} TODO: the item that surfaces the target misconception, and
        that the debrief opens on.
        \answerspace{1.8cm}{}
  \item TODO: the follow-through in context.
        \answerspace{1.8cm}{}
\end{enumerate}
\end{scenariobox}

\vspace{0.10in}

% ── DEBRIEF (8 min, whole class — filled together after the groups report out) ─
% Three synthesis prompts, short. Item 1 is the sentence you want on every desk.
\boxguard[18]
\begin{reflectionbox}
\small
\textbf{Debrief --- we fill this in together.}
\begin{enumerate}[leftmargin=1.9em, itemsep=6pt]
  \item TODO: state the lesson's key idea in one sentence.
        \answerspace{1.0cm}{}
  \item TODO: generate an example \emph{not} on this page --- where transfer shows up.
        \answerspace{1.2cm}{}
  \item TODO: the ``why does this matter'' line, in their own words.
        \answerspace{1.0cm}{}
\end{enumerate}
\end{reflectionbox}

\end{document}
